% Section 5: The arithmetic--process dictionary and the equivalence.

\subsection{What an equivalence must consist of}

An equivalence between arithmetic expression geometry and process geometry is not a
single theorem.  It consists of four layers, and each layer has its own strength:

\begin{enumerate}
\item \emph{Objects}: an expression tree corresponds to a program body; a marked
  spinal history corresponds to a schedule; a chirality word corresponds to a slot
  order.
\item \emph{Orders}: the dependency order of Paper~I corresponds to the construction
  order of the process; the chronological order of a sequential history corresponds
  to the unique linear extension of a single-spine program.
\item \emph{Carrier}: the six-state carrier of Paper~V corresponds to the iota frame
  of the process corpus, plane by plane.
\item \emph{Structure}: the arithmetic structures --- one-hole contexts, projective
  evaluation, charge --- correspond to process structures (graft traces, observations,
  cut linkage), or they do not, and the failure is recorded.
\end{enumerate}

The paper states the dictionary at the first three layers and treats the fourth as a
programme: the layers already established are listed with status, the layers that are
only proposed are marked, and the layers that fail are recorded as boundaries rather
than patched.

\subsection{The dictionary}

\begin{center}
\footnotesize
\begin{longtable}{@{}>{\raggedright\arraybackslash}p{0.32\linewidth}>{\raggedright\arraybackslash}p{0.38\linewidth}>{\raggedright\arraybackslash}p{0.24\linewidth}@{}}
\toprule
Arithmetic expression geometry & Process geometry & Status \\
\midrule
expression tree (Paper~I) & program body with ordered slots
  & \texttt{DEFINITION} \\
free marked spinal history (Paper~I) & schedule of a single-spine program
  & Proposition~\ref{prop:sequential-leaf}, \texttt{PROVED} \\
chirality word (Paper~I) & slot order induced by the arity-one carrier
  & \texttt{STRUCTURAL PROPOSAL} \\
spine and left/right combs (Paper~I) & single-spine and path cases of a slot order
  & \texttt{STRUCTURAL PROPOSAL} \\
dependency poset (Paper~I) & construction order (Definition~\ref{def:construction-order})
  & \texttt{DEFINITION} (tree bodies) \\
corner order and linear extensions & schedules and their inversions
  & Theorem~\ref{thm:canonical-order}, \texttt{PROVED} \\
evaluation order of a sequential history & unique linear extension of a chain
  & \texttt{PROVED} (Proposition~\ref{prop:sequential-leaf}) \\
six-state carrier (Paper~V) & iota frame of the process corpus
  & \texttt{CORRESPONDENCE} (note~40) \\
chirality sign law (Paper~V) & sign action of $J$ on the planes
  & \texttt{CORRESPONDENCE} \\
carousel class four-cycle & non-canonical alignment with $J$'s four-cycle plus a shift
  & \texttt{STRUCTURAL PROPOSAL} + \texttt{CORRESPONDENCE} \\
carousel count shadow & involution $M^{2}=\mathrm{id}$
  & \texttt{CORRESPONDENCE} (source) + calibration \\
one-hole contexts and projective evaluation & observations compiled from a program
  & \texttt{OPEN PROBLEM} \\
accumulative commutative charge & cut linkage / frontier data
  & \texttt{OPEN PROBLEM} \\
\bottomrule
\end{longtable}
\end{center}

\begin{remark}[Reading the table]\label{rem:reading-table}
The table is a typed dictionary, not a theorem.  Each row asserts that a construction
on one side has a counterpart on the other, at the strength named.  The rows marked
\texttt{OPEN} are directions, not identifications: no claim of this paper depends on
them, and the equivalence statements below do not use them.
\end{remark}

\subsection{The equivalence statements}

\begin{theorem}[Sequential recovery]\label{thm:sequential-recovery}
The single-spine fragment of the dictionary is an equivalence on objects: the map
sending a \emph{free} marked spinal history to its chain program and back is a
bijection, compatible with concatenation of one-hole contexts, between free marked spinal
histories with their chronological order and open finite programs with one hole and a
chain body whose unary contexts carry the Paper~I signature, with the slot convention
declared as data.  A bounded history keeps its initial mark as an extra datum, outside the
bijection.  Under it, the
chirality word is the slot sequence, the unique schedule is the chronological order,
and the spine is the body chain.
\end{theorem}

\begin{proof}
Both directions are the explicit construction of
Proposition~\ref{prop:sequential-leaf}; compatibility is immediate because both sides
compose by concatenation of contexts along the slot sequence, and the slot convention
is the only choice involved.  No statement is made about the branching case.
\end{proof}

\begin{conjecture}[General process--expression equivalence]\label{conj:general-equivalence}
There is a natural assignment from plane binary trees equipped with the corner order
and dependency order to open finite programs equipped with a planar schedule
structure --- natural meaning compatible with subtree substitution on the two sides ---
which restricts to Theorem~\ref{thm:sequential-recovery}
on the single-spine fragment and sends the canonical evaluation order to the unique
planar schedule closest to the corner order.
\end{conjecture}

The conjecture is deliberately narrow: it asks for an assignment on the planar layer, not
for an identification of arbitrary programs with expressions.  What is presently
known toward it is: the sequential fragment (Theorem~\ref{thm:sequential-recovery});
the canonical order on all plane trees (Theorem~\ref{thm:canonical-order}); the
carrier identification (\texttt{CLASSICAL}); and the screw calibration
(\S\ref{sec:carousel}).  What is missing is a general construction of the assignment and
a proof that it preserves the schedule structure; both are recorded as open.

\begin{proposal}[What the equivalence is for]\label{prop:equivalence-role}
The purpose of the equivalence is to let one theory's established results be reused
in the other without collapsing their objects.  Concretely: an order-theoretic result
about schedules becomes a statement about histories, and a geometric statement about
the carrier becomes a statement about the process's chirality.  It is not a licence to
identify values with programs, schedules with processes, or charges with frontier
data.
\end{proposal}

\subsection{Finite tests the equivalence must pass}

Three finite tests are stated so that later work can refute the conjecture rather
than illustrate it.

\begin{enumerate}
\item \emph{Tree census}: for every plane binary tree with at most six internal
  vertices (seven leaves), the dictionary
  image of the canonical order is the canonical schedule of the image program, and the
  inversion count is preserved.  The tree side is already calibrated; the program side
  needs an executed fixture.
\item \emph{Schedule count}: the number of schedules of the image of a tree equals the
  number of linear extensions of its dependency order, and the canonical schedule is
  the image of the canonical extension.  The count is the tree's hook-length product;
  a mismatch is a refutation.
\item \emph{Carrier check}: the dictionary's carrier map sends the sign law of
  Paper~V to the sign action of $J$, including the transposition reversal; a mismatch
  at any of the three planes is a refutation.
\end{enumerate}

None of the three tests is currently stored as an executed fixture on the program
side; they are stated here as acceptance tests for the programme and are registered in
\S\ref{sec:boundaries}.  This is the honest status: the equivalence is proved on the
sequential fragment, supported by the carrier identification and the carousel
calibration, and conjectural in general.
